%% ═══════════════════════════════════════════════════════════════════════════════
\section{Attack Constraint Variants}
\label{sec:modes}
%% ═══════════════════════════════════════════════════════════════════════════════

In this section, we describe the two variants for encoding constraint T3 (``$N_2$ is fooled by $\xp$'') in the transfer MIP. The choice of variant determines whether the adversary targets a specific class or any misclassification.

\subsection{Targeted Attack}

In the targeted variant, the adversary forces $N_2$ to predict a specific target class $t$. Formally, the constraint requires:
\begin{equation}
  \label{eq:ctarget}
  \hat{y}_t\bigl(N_2(\xp)\bigr) \;\ge\; \hat{y}_k\bigl(N_2(\xp)\bigr)
  \quad \forall\, k \neq t,
\end{equation}
equivalently $\conf{N_2}{\xp}{t} \ge 0$. This is the more restrictive variant: the adversary must not only cause a misclassification but must steer the network toward a particular wrong class.

\subsection{Untargeted Attack}

In the untargeted variant, the adversary need only reduce $N_2$'s confidence in the true class $c$ to negative, causing any misclassification:
\begin{equation}
  \label{eq:ctag}
  \conf{N_2}{\xp}{c} \;\le\; 0.
\end{equation}

\subsection{Comparison}

Table~\ref{tab:modes} compares the two variants. The targeted variant has a smaller feasible region (the adversary must achieve a specific misclassification), leading to a larger or equal $\dD$ and thus a stronger transfer guarantee. The untargeted variant has a larger feasible region, yielding a smaller or equal $\dD$ and a weaker but more general guarantee.

\begin{table}[ht]
  \centering
  \caption{Comparison of targeted and untargeted attack constraints.}
  \label{tab:modes}
  \begin{tabular}{lll}
    \toprule
    \textbf{Property} & \textbf{Targeted} & \textbf{Untargeted} \\
    \midrule
    Attack type
      & $c \to t$ (specific class) & $c \to \text{any}$ \\
    Constraint on $N_2(\xp)$
      & $\hat{y}_t \ge \hat{y}_k,\;\forall k\neq t$
      & $\hat{y}_c < \hat{y}_k$ for some $k$ \\
    Feasible region    & Smaller (harder to satisfy) & Larger (easier) \\
    Resulting $\dD$    & Larger or equal & Smaller or equal \\
    Inference strength & Stronger guarantee & Weaker guarantee \\
    Typical use        & Specific adversary model & General robustness \\
    \bottomrule
  \end{tabular}
\end{table}

